% titel: Blatt 9, A2: Gefrierpunkt- und Siedepunkterhöhung (kolligativ)
% skript: Mortimer Kap. 13 (kolligative Eigenschaften)

% Teil I: Werkzeuge, kein Lösungsweg.

Ein gelöster nicht-flüchtiger Stoff erniedrigt den Gefrierpunkt und erhöht
den Siedepunkt des Lösungsmittels.

\begin{werkzeugkasten}[Kolligative Effekte]
Gefrierpunkterniedrigung und Siedepunkterhöhung hängen nur von der
Molalität der Teilchen ab (nicht von der Art des Stoffes):
\[
  \Delta T_f = i \cdot E_f \cdot b,\qquad
  \Delta T_b = i \cdot E_b \cdot b.
\]
\begin{itemize}
  \item $E_f=1.86\,K\,kg\,mol^{-1}$, $E_b=0.513\,K\,kg\,mol^{-1}$ für Wasser.
  \item $b$ = Molalität (mol gelöster Substanz pro kg \textbf{Lösungsmittel}).
  \item $i$ = Van't-Hoff-Faktor (Zahl der Teilchen je Formeleinheit in
        Lösung; bei Nicht-Elektrolyten $i=1$).
\end{itemize}
Gefrierpunkt $T_f = 0\,$°C $-$ $\Delta T_f$; Siedepunkt $T_b = 100\,$°C $+$
$\Delta T_b$.
\textit{Skript:} Mortimer Kap. 13 (kolligative Eigenschaften).
\end{werkzeugkasten}

\begin{achtung}[Darauf achten]
\begin{itemize}
  \item Glycerin dissoziiert nicht ($i=1$), $\Rightarrow$ einfach
        $b$ aus Masse/Molmasse/$\mathrm{kg}$ Lösungsmittel.
  \item Lösungsmittelmasse in kg einsetzen (120.0 g $=$ 0.120 kg).
  \item Der Gefrierpunkt ist niedriger, der Siedepunkt höher als beim reinen
        Wasser --- also $T_f = 0-\Delta T_f$ und $T_b = 100+\Delta T_b$.
\end{itemize}
\end{achtung}
